\documentclass[10pt,a4paper]{article}
\usepackage[T1]{fontenc}
\usepackage{lmodern}
\usepackage[left=24mm,right=24mm,top=20mm,bottom=18mm]{geometry}
\usepackage{amsmath,amssymb,microtype,booktabs,fancyhdr}
\usepackage[hidelinks]{hyperref}
\hypersetup{pdftitle={RoPE attention is an exact forward-pass gradient step with softmax intact.},
  pdfauthor={Julie Huang, Maggie Chlon, Leon Chlon},
  pdfsubject={Exact query-conditioned gradient steps and tokenwise residual decomposition}}
\setlength{\parindent}{0pt}
\setlength{\parskip}{10pt}
\setlength{\abovedisplayskip}{12pt}
\setlength{\belowdisplayskip}{13pt}
\setlength{\abovedisplayshortskip}{12pt}
\setlength{\belowdisplayshortskip}{13pt}
\setlength{\headheight}{12pt}
\pagestyle{fancy}
\fancyhf{}
\renewcommand{\headrulewidth}{0pt}
\fancyfoot[L]{\footnotesize Exact RoPE--softmax construction}
\fancyfoot[R]{\footnotesize\thepage\,/\,4}
\renewcommand{\arraystretch}{1.35}
\newcommand{\R}{\mathbb R}
\newcommand{\T}{^{\!\top}}
\newcommand{\calI}{\mathcal I}
\newcommand{\ddM}{\Delta M}
\newcommand{\norm}[1]{\left\lVert#1\right\rVert}
\begin{document}
\fontsize{10}{15.8}\selectfont
\begin{center}
{\fontsize{16}{19}\selectfont\bfseries RoPE attention is an exact forward-pass\\gradient step with softmax intact.}\par
\vspace{6pt}
{\fontsize{11.5}{15}\selectfont Julie Huang\textsuperscript{1}, Maggie Chlon\textsuperscript{1}, Leon Chlon\textsuperscript{1,2,*}}\par
{\fontsize{10.5}{14}\selectfont \textsuperscript{1}Hassana Labs;\quad\textsuperscript{2}Department of Information Engineering, Oxford University}\par
{\fontsize{10.5}{14}\selectfont \textsuperscript{*}Correspondence: \href{mailto:leochlon@robots.ox.ac.uk}{leochlon@robots.ox.ac.uk}}
\end{center}
\vspace{-3pt}
\fontsize{10}{15.8}\selectfont
\setlength{\parskip}{10pt}

\textbf{Contributions: exact inference-time gradient identity and tokenwise decomposition}

We express RoPE--softmax attention as an exact gradient step for arbitrary projection
weights and affine biases. The construction places the head's forward pass in
effective-weight coordinates and resolves its query-dependent write into cached-token
contributions. Full softmax and rotary position structure remain intact throughout.

\textbf{An explicit forward-pass gradient step.}
Use the augmented input $u_i$ and projections of section 3; write
$L_i=\widehat W_qR(p_i)\T/\tau$, $a_{ij}=L_i\widetilde k_j\T$, and
$c_{ij}=\rho(s_{ij})/Z_i>0$. For $B\in\R^{(m+1)\times r}$ define
\begin{equation}
\boxed{\begin{aligned}
\mathcal E_i(B)&=\frac12\sum_{j\in\calI_i}c_{ij}
 \norm{a_{ij}\T B-(v_j-\mu_i)}_2^2,\\
B_i^{(1)}&=0-\nabla_B\mathcal E_i(0)
 =\sum_{j\in\calI_i}c_{ij}a_{ij}(v_j-\mu_i)=\ddM_i,
\qquad y_i=\mu_i+u_i\T B_i^{(1)}.
\end{aligned}}
\label{eq:regression}
\end{equation}
Differentiation acts on $B$, with features, targets, and coefficients held fixed.
Each query specifies the positive regression weights and reads out the resulting
unit step from zero.

\textbf{Exact tokenwise residual.}
At a fixed cache and query position, write $a_j=a_{ij}$ and choose a reference
query $0$. The adaptive write equals the reference write plus
\begin{equation}
\boxed{\displaystyle
e_i=u_i\T(\ddM_i-\ddM_0)
=\sum_{j\in\calI}(c_{ij}-c_{0j})(u_i\T a_j)(v_j-\mu).}
\label{eq:tokenerror}
\end{equation}
Each summand identifies a cached token's contribution to the correction.
Thus $y_i=\mu+u_i\T\ddM_0+e_i$, and $e_iW_o$ is the exact projected correction.

\textbf{RoPE as a bank of positional derivatives.}
For fixed content $z_0$, $z(p)=R(p)z_0=e^{pA}z_0$ satisfies
$\partial_p^n z(p)=A^nz(p)$ for $n\geq1$, and
$z(p+\delta)-z(p)=(e^{\delta A}-I)z(p)$.
These derivatives and finite differences are directly evaluable on frozen features.
Our framework connects this derivative-operator representation to the exact
full-softmax effective matrix. RollPE relates a special RoPE spectrum to Fourier
differentiation [6].

\textbf{Shared cache, query-conditioned weights.}
At a fixed cache and position, $\ddM_i=LT_i$: $T_i\in\R^{d\times r}$ combines
cached outer products with positive coefficients $c_{ij}$, and
$\operatorname{rank}(\ddM_iW_o)\leq d$ in the original gauge. The mean is shared.
The matrix gives a secant representation; differentiating the input and the
coefficients gives the query Jacobian.

\textbf{Equivalent matrices, one output map.}
For any fixed $g\in\R^{m+1}$, set $\sigma_i=u_i\T g$. Then
\begin{equation}
\begin{aligned}
a_{ij}^{g}&=a_{ij}-g,\qquad s_{ij}^{g}=s_{ij}-\sigma_i,
\qquad c_{ij}^{g}=\frac{e^{\sigma_i}\rho(s_{ij}-\sigma_i)}{Z_i},\\
\ddM_i^{g}&=\sum_j c_{ij}^{g}a_{ij}^{g}(v_j-\mu_i),
\qquad y_i=\mu_i+u_i\T\ddM_i^{g}.
\end{aligned}
\label{eq:gauge}
\end{equation}
These transformations preserve the probabilities and output. Key shifts give the
subfamily $g=L_i b$, of dimension at most $d$; general $g$ can add one direction.
We report matrices in the original-score gauge and evaluate transfer in output coordinates.

\textbf{Connection to linear-attention gradient constructions.}
Von Oswald's Proposition 1 [1] realizes a least-squares gradient step with linear
attention; related fast-weight [2] and implicit-gradient [3] constructions use
linear attention or linearized softmax. Setting $R=I$ and taking $c_{ij}\to1/N_i$,
with centered regression tokens, his projection/readout choices, and matched scale,
recovers his step. Equation~\eqref{eq:regression} supplies the exact full-softmax coefficients.

\textbf{Pretrained realization.}
On Qwen2.5-0.5B, layer 23, all 240 context/head reconstruction rows pass.
The combined prefix and full-attention writes reconstruct at float64 MSE
$9.78\times10^{-30}$ and $9.62\times10^{-30}$, respectively.

\medskip

\fontsize{10}{15.8}\selectfont
\setlength{\parskip}{10pt}
\textbf{Mathematical construction}

\textbf{1. RoPE and the score in the original weights.}
Let $x_i\in\R^m$ be a column input to one head, with
$W_q,W_k\in\R^{m\times d}$ and $W_v\in\R^{m\times r}$.
Projected vectors are rows: $q_i=x_i\T W_q$, $k_j=x_j\T W_k$,
and $v_j=x_j\T W_v$. Set $\tau=\sqrt d$.
For standard RoPE with fixed frequencies [4],
\begin{equation}
R(p)=e^{pA},\qquad
A=\operatorname{diag}_{\ell=1}^{d/2}
\begin{pmatrix}0&-\omega_\ell\\\omega_\ell&0\end{pmatrix},
\qquad R(p_i)\T R(p_j)=R(p_j-p_i).
\label{eq:rope}
\end{equation}
With
$\widetilde q_i=q_iR(p_i)\T$ and $\widetilde k_j=k_jR(p_j)\T$, define
\begin{equation}
a_{ij}=\frac{W_qR(p_j-p_i)W_k\T x_j}{\tau}\in\R^m,
\qquad
s_{ij}=\frac{\widetilde q_i\widetilde k_j\T}{\tau}=x_i\T a_{ij}.
\label{eq:score}
\end{equation}
\textbf{2. Exact factorization through softmax.}
For deterministic attention, let $\calI_i$ be a nonempty set of permitted keys;
causal masking selects this set. With the scores in \eqref{eq:score}, write
$N_i=|\calI_i|$ and
\begin{equation}
Z_i=\sum_{j\in\calI_i}e^{s_{ij}},\qquad
\mu_i=\frac1{N_i}\sum_{j\in\calI_i}v_j,\qquad
y_i=\frac1{Z_i}\sum_{j\in\calI_i}e^{s_{ij}}v_j.
\label{eq:softmax}
\end{equation}
Here $\mu_i$ is the \emph{uniform arithmetic mean} of permitted values. The divided difference
\begin{equation}
\rho(s)=\int_0^1 e^{ts}\,dt,\quad
\rho(s)=\frac{e^s-1}{s}\ (s\ne0),\quad
\rho(0)=1,\quad e^s-1=s\rho(s).
\label{eq:rho}
\end{equation}
gives, for every query,
\begin{equation}
\boxed{\begin{aligned}
 y_i&=\mu_i+x_i\T\ddM_i,\\
 \ddM_i&=\frac1{\tau Z_i}\sum_{j\in\calI_i}
 \rho(s_{ij})\big[W_qR(p_j-p_i)W_k\T x_j\big]
 \big[x_j\T W_v-\mu_i\big].
\end{aligned}}
\label{eq:main}
\end{equation}
\emph{Proof.} The centered values sum to zero. With all sums over $\calI_i$,
\begin{equation*}
y_i-\mu_i
=\frac{\sum_j(e^{s_{ij}}-1)(v_j-\mu_i)}{Z_i}
=\frac{\sum_j s_{ij}\rho(s_{ij})(v_j-\mu_i)}{Z_i}
=x_i\T\ddM_i.\qquad\square
\end{equation*}
Setting $R=I$ gives the position-free construction; RoPE supplies relative-position structure.

\textbf{Convention.}
Reported matrices use the \emph{original-score gauge} of \eqref{eq:main}; gauge
changes transform both features and coefficients, as in \eqref{eq:gauge}.

\textbf{Transformed scores, additive biases, and QK normalization.}
For pointwise $\phi$ with $\phi(0)=0$, use $Z=\sum_j e^{\phi(s_{ij})}$ and
$\rho_\phi(s)=(e^{\phi(s)}-1)/s$ for $s\ne0$; at zero, any finite value works
(continuously, $\phi'(0)$ when it exists). Sign-preserving soft-caps preserve
nonnegative regression weights. Additive input-independent biases, including ALiBi,
are absorbed by $(x_i\T,1)(a_{ij}\T,b_{ij})\T$.
For nonlinear QK normalization (Qwen3 [5]), use the normalized, rotated query:
\begin{equation}
y_i=\mu_i+\widetilde q_i S_i,\qquad
S_i=\frac1\tau\sum_j\frac{\rho(s_{ij})}{Z_i}\widetilde k_j\T(v_j-\mu_i).
\label{eq:qknorm}
\end{equation}

\textbf{Finite-precision evaluation in the declared gauge.}
With $M_i=\max_j s_{ij}$,
$\overline Z_i=\sum_j e^{s_{ij}-M_i}$ and $\ell_i=M_i+\log\overline Z_i$,
\begin{equation}
c_{ij}=\frac{P_{ij}-e^{-M_i}/\overline Z_i}{s_{ij}}
=\frac{P_{ij}-e^{-\ell_i}}{s_{ij}}\ (s_{ij}\ne0),\qquad
c_{ij}=e^{-\ell_i}\ (s_{ij}=0).
\label{eq:stable}
\end{equation}
Evaluate near-zero scores with $e^{-\ell_i}\operatorname{expm1}(s_{ij})/s_{ij}$
and use log-domain arithmetic elsewhere. Numerical scaling preserves the original gauge.

\medskip

\fontsize{10}{15.8}\selectfont
\setlength{\parskip}{10pt}
\textbf{3. Effective residual weights and affine projections.}
With $W_o\in\R^{r\times m}$ and residual input $x_i$, equation~\eqref{eq:main} gives
\begin{equation}
x_i\T+y_iW_o=x_i\T(I_m+\ddM_iW_o)+\mu_iW_o.
\label{eq:residual}
\end{equation}
For projection biases, use $u_j=(x_j\T,1)\T$ and
$\widehat W=\bigl[W\T\;b\bigr]\T$; the proof is unchanged.
For pre-normalization, $x_i=\operatorname{Norm}(h_i)$ and the residual is
$h_i\T+(\mu_i+u_i\T\ddM_i^{\mathrm{aug}})W_o$.
Add output bias separately and sum projected heads, each with its assigned K/V bank.

\textbf{4. Exact reference decomposition and matrix transfer.}
Fix a context bank and query position. For anchor $0$, set
$y_i^{\mathrm{fr}}=\mu+x_i\T\ddM_0$. Then
\begin{equation}
\boxed{\begin{aligned}
y_i-y_i^{\mathrm{fr}}&=x_i\T(\ddM_i-\ddM_0),\\
\norm{y_i-y_i^{\mathrm{fr}}}_2&\leq\norm{x_i}_2\,\norm{\ddM_i-\ddM_0}_{\mathrm{op}}.
\end{aligned}}
\label{eq:freeze}
\end{equation}
The contraction in \eqref{eq:freeze} is the exact transfer residual.
With $L=W_qR(p)\T/\tau$ and $c_{ij}=\rho(s_{ij})/Z_i$,
\begin{equation}
\ddM_i-\ddM_0=L\sum_{j\in\calI}(c_{ij}-c_{0j})\widetilde k_j\T(v_j-\mu).
\label{eq:refresh}
\end{equation}
These equations hold in augmented coordinates. Reference write plus coefficient
refresh gives the exact readout. Freezing omits $x_i\T(\ddM_i-\ddM_0)$;
equation~\eqref{eq:tokenerror} resolves that residual by cached token.

\textbf{5. Fixed-prefix and full-attention readouts.}
A fixed-prefix probe excludes the query's own key/value to keep one common bank.
If full attention adds a self value $v_i^{\mathrm{self}}$ with softmax mass $\pi_i$, then
\begin{equation}
y_i^{\mathrm{full}}=(1-\pi_i)y_i^{\mathrm{prefix}}+\pi_i v_i^{\mathrm{self}}.
\label{eq:full}
\end{equation}
The full-bank matrix uses its own scores and mean
$(N\mu+v_i^{\mathrm{self}})/(N+1)$. A changing bank adds $\mu_i-\mu_0$ to the
error when the offset is also frozen.

\textbf{6. Pretrained validation and frozen reuse.}
Qwen2.5-0.5B, layer 23 (zero-indexed): $C=16$ prefixes of 511 tokens,
all 14 heads, $Q=32$ continuations at one position. Anchor $a=0$ is each article's
observed next token; the others are the model's 31 highest-ranked distinct nonspecial
continuations. Each reads the same prefix-only bank.

\textbf{Exact reconstruction.}
All 240 context/head reconstruction rows pass. Comparison with the native FP32
write gives MSE $7.01\times10^{-14}$ and maximum absolute error $1.4\times10^{-5}$.

\textbf{Frozen reuse.}
For each prefix, reuse the anchor's matrix across its other queries.
Let
$o_{ci}=\sum_h y_{cih}W_{oh}\in\R^D$ denote the projected attention write.
\begin{equation}
E(f)=\frac1C\sum_{c=1}^C\frac1{(Q-1)D}\sum_{i\ne a}
\norm{f_{ci}-o_{ci}}_2^2,\qquad
\mathrm{rRMSE}(f)=\sqrt{\frac{E(f)}{E(o_{ca})}}.
\label{eq:metrics}
\end{equation}
Pool squared errors with equal context weights and take ratios after pooling.
All-anchor evaluation also averages over $a=0,\ldots,Q-1$ before taking ratios.
Sum projected heads before squaring. Head-level baselines are
\begin{equation}
y_i^{\mu}=\mu,\qquad
y_i^{\mathrm{lin}}=\mu+\frac1N\sum_j s_{ij}(v_j-\mu).
\label{eq:baselines}
\end{equation}
The latter is the first-order softmax expansion at zero logits with unit scale.
Both undergo the same output projection and aggregation as the exact head.

Pooled MSE of the combined projected write:
\begin{equation*}
\begin{array}{lrr}
\text{Predictor} & \text{Anchor }0 & \text{All anchors}\\
\text{Constant anchor output} & 0.047787 & 0.045672\\
\text{Cache mean }(\ddM=0) & 0.124573 & 0.124513\\
\text{Frozen anchor matrix} & 0.324970 & 0.356763\\
\text{Linearized softmax} & 0.161842 & 0.162068
\end{array}
\end{equation*}
Target RMS is $0.670545$ for anchor $0$ and $0.670824$ across all anchors.
Frozen rRMSE is $2.608$ against anchor $0$ and $1.615$ against the cache mean;
the all-anchor rRMSE is $2.795$. In the all-anchor evaluation, thirteen individual
heads have rRMSE $0.728$--$0.933$ against their constant-anchor baselines.

\medskip
\textbf{References}

{\footnotesize [1] J. von Oswald et al., \emph{Transformers Learn In-Context by Gradient
Descent}, \href{https://proceedings.mlr.press/v202/von-oswald23a.html}{ICML 2023}.}

{\footnotesize [2] I. Schlag, K. Irie, J. Schmidhuber, \emph{Linear Transformers Are Secretly Fast
Weight Programmers}, \href{https://proceedings.mlr.press/v139/schlag21a.html}{ICML 2021}.}

{\footnotesize [3] D. Dai et al., \emph{Why Can GPT Learn In-Context? Language Models Implicitly
Perform Gradient Descent as Meta-Optimizers}, \href{https://aclanthology.org/2023.findings-acl.247/}{Findings of ACL 2023}.}

{\footnotesize [4] J. Su et al., \emph{RoFormer: Enhanced Transformer with Rotary Position
Embedding}, \href{https://arxiv.org/abs/2104.09864}{2021}.}

{\footnotesize [5] Qwen team and Hugging Face, \href{https://github.com/huggingface/transformers/blob/main/src/transformers/models/qwen3/modeling_qwen3.py}{Qwen3 attention implementation}, Q/K RMS normalization.}

{\footnotesize [6] C. van de Geijn et al., \emph{Do traveling waves make good positional encodings?},
\href{https://arxiv.org/html/2511.11668v1}{2025}, appendices E--F.}

\end{document}
